\documentclass[11pt,a4paper]{article}
\usepackage[utf8]{inputenc}
\usepackage{amsmath,amssymb,amsthm}
\usepackage{algorithm}
\usepackage{algpseudocode}
\usepackage{booktabs}
\usepackage{hyperref}
\usepackage{geometry}
\geometry{margin=1in}

\title{Additive and Multiplicative Recurrent Attention: Bahdanau Alignment, Luong Formulation, and Differentiable Memory Synthesis}
\author{Team Scaibu \\ \small Email: \texttt{jaydeep.9664920749@gmail.com} \\ \small Architecture and Core Engine Team}
\date{\today}

\newtheorem{theorem}{Theorem}
\newtheorem{lemma}[theorem]{Lemma}
\newtheorem{proposition}[theorem]{Proposition}
\theoremstyle{definition}
\newtheorem{definition}{Definition}
\newtheorem{remark}{Remark}

\begin{document}

\maketitle

\begin{abstract}
Fixed-vector sequence-to-sequence architectures suffer from unavoidable information degradation when encoding long source sequences. Neural attention mechanisms overcome this vector bottleneck by providing decoders with direct, differentiable soft-alignment access across all intermediate encoder hidden states. This document presents the mathematical formulations of Additive (Bahdanau) and Multiplicative (Luong) attention scoring functions, establishes gradient propagation properties through convex attention mixtures, details algorithmic pseudocode, and provides a verified step-by-step numerical calculation.
\end{abstract}

\section{Notation and Mathematical Preliminaries}

Let $\mathbf{h} = (\mathbf{h}_1, \dots, \mathbf{h}_{T_x})$ denote the sequence of encoder annotation vectors with $\mathbf{h}_i \in \mathbb{R}^{d_h}$, and let $\mathbf{s}_t \in \mathbb{R}^{d_s}$ denote the decoder query state vector at decoding time step $t$.

\begin{table}[h]
\centering
\caption{Mathematical Notation for Recurrent Attention Mechanisms}
\begin{tabular}{lll}
\toprule
\textbf{Symbol} & \textbf{Domain} & \textbf{Semantic Description} \\
\midrule
$e_{t, i}$ & $\mathbb{R}$ & Continuous alignment energy score between decoder $t$ and encoder $i$ \\
$\boldsymbol{\alpha}_t$ & $\Delta^{T_x - 1}$ & Softmax attention alignment probability distribution on the simplex \\
$\mathbf{c}_t$ & $\mathbb{R}^{d_h}$ & Synthesized dynamic context vector at decoding step $t$ \\
$W_a, U_a, \mathbf{v}_a$ & Parameters & Bahdanau additive projection matrices and alignment vector \\
$W_{\text{luong}}$ & $\mathbb{R}^{d_s \times d_h}$ & Luong general multiplicative bilinear alignment matrix \\
$\tilde{\mathbf{s}}_t$ & $\mathbb{R}^{d_s}$ & Attentional hidden state vector \\
\bottomrule
\end{tabular}
\end{table}

\section{Problem Definition and Core Concepts}

\subsection{3-Level Explanation}
\begin{enumerate}
    \item \textbf{Intuitive (High School level):} Instead of memorizing a whole book before translating, the translator keeps the source book open. When writing each word, they glance at the exact original sentence that matters most, highlighting key words with soft focus.
    \item \textbf{Engineering Level:} Attention dynamically computes a probability distribution $\boldsymbol{\alpha}_t$ over all $T_x$ source hidden states. The resulting context vector $\mathbf{c}_t = \sum_i \alpha_{t, i} \mathbf{h}_i$ injects exact source information into the decoder prediction head, eliminating encoder compression bottlenecks.
    \item \textbf{Mathematical Level:} Given decoder state $\mathbf{s}$ and encoder states $\{\mathbf{h}_i\}$, the attention operator maps $(\mathbf{s}, \{\mathbf{h}_i\}) \mapsto \mathbf{c}$ via:
    \begin{equation}
    e_i = \text{Score}(\mathbf{s}, \mathbf{h}_i), \quad \alpha_i = \frac{\exp(e_i)}{\sum_j \exp(e_j)}, \quad \mathbf{c} = \sum_{i=1}^{T_x} \alpha_i \mathbf{h}_i
    \end{equation}
    where $\text{Score}_{\text{Bahdanau}} = \mathbf{v}_a^T \tanh(W_a \mathbf{s} + U_a \mathbf{h}_i)$ and $\text{Score}_{\text{Luong}} = \mathbf{s}^T W_a \mathbf{h}_i$.
\end{enumerate}

\subsection{5-Step Algorithmic Framework}
\begin{enumerate}
    \item \textbf{Linear State Projection:} Project decoder state and encoder representations into shared alignment space $\mathbb{R}^{d_a}$.
    \item \textbf{Energy Score Computation:} Evaluate non-linear additive or bilinear dot-product scores $e_{t, i}$.
    \item \textbf{Softmax Normalization:} Exponentiate and normalize energy scores into probability distribution $\boldsymbol{\alpha}_t$.
    \item \textbf{Context Vector Synthesis:} Compute convex combination $\mathbf{c}_t = \sum_i \alpha_{t, i} \mathbf{h}_i$.
    \item \textbf{Attentional State Combination:} Fuse context $\mathbf{c}_t$ and decoder state $\mathbf{s}_t$ through non-linear emission head $\tilde{\mathbf{s}}_t = \tanh(W_c [\mathbf{c}_t; \mathbf{s}_t])$.
\end{enumerate}

\section{Algorithm Specification}

\begin{algorithm}[H]
\caption{Bahdanau Additive Attention Context Synthesis}
\label{alg:bahdanau}
\begin{algorithmic}[1]
\Require Decoder query $\mathbf{s} \in \mathbb{R}^{d_s}$, encoder keys $\mathbf{H} \in \mathbb{R}^{T_x \times d_h}$, parameters $W_a \in \mathbb{R}^{d_a \times d_s}, U_a \in \mathbb{R}^{d_a \times d_h}, \mathbf{v}_a \in \mathbb{R}^{d_a}, \mathbf{b}_a \in \mathbb{R}^{d_a}$.
\Ensure Context vector $\mathbf{c} \in \mathbb{R}^{d_h}$, attention weights $\boldsymbol{\alpha} \in \mathbb{R}^{T_x}$.
\State $\mathbf{w}_s \gets W_a \mathbf{s} + \mathbf{b}_a$
\For{$i = 1$ \textbf{to} $T_x$}
    \State $\mathbf{u}_i \gets U_a \mathbf{H}[i]$
    \State $e_i \gets \mathbf{v}_a^T \tanh(\mathbf{w}_s + \mathbf{u}_i)$
\EndFor
\State $e_{\max} \gets \max_i e_i$
\State $S \gets \sum_{i=1}^{T_x} \exp(e_i - e_{\max})$
\For{$i = 1$ \textbf{to} $T_x$}
    \State $\alpha_i \gets \exp(e_i - e_{\max}) / S$
\EndFor
\State $\mathbf{c} \gets \sum_{i=1}^{T_x} \alpha_i \mathbf{H}[i]$
\State \Return $\mathbf{c}, \boldsymbol{\alpha}$
\end{algorithmic}
\end{algorithm}

\section{Theoretical Guarantees and Gradient Backpropagation}

\begin{theorem}[Direct Gradient Flow Path]
\label{thm:attention_gradient}
Under attention context synthesis $\mathbf{c}_t = \sum_{i=1}^{T_x} \alpha_{t, i} \mathbf{h}_i$, the gradient of the loss with respect to source annotation $\mathbf{h}_i$ contains a direct path:
\begin{equation}
\frac{\partial \mathcal{L}}{\partial \mathbf{h}_i} = \alpha_{t, i} \frac{\partial \mathcal{L}}{\partial \mathbf{c}_t} + \sum_{j=1}^{T_x} \frac{\partial \mathcal{L}}{\partial \alpha_{t, j}} \frac{\partial \alpha_{t, j}}{\partial \mathbf{h}_i}
\end{equation}
guaranteeing $\mathcal{O}(1)$ temporal path length between arbitrary source positions and target output emissions, bypassing vanishing gradients.
\end{theorem}

\begin{proof}
By total derivative of loss $\mathcal{L}(\mathbf{c}_t(\mathbf{h}_1, \dots, \mathbf{h}_{T_x}), \boldsymbol{\alpha}_t(\mathbf{h}_1, \dots, \mathbf{h}_{T_x}))$ with respect to $\mathbf{h}_i$:
$\frac{\partial \mathcal{L}}{\partial \mathbf{h}_i} = \frac{\partial \mathcal{L}}{\partial \mathbf{c}_t} \frac{\partial \mathbf{c}_t}{\partial \mathbf{h}_i} + \sum_j \frac{\partial \mathcal{L}}{\partial \alpha_{t, j}} \frac{\partial \alpha_{t, j}}{\partial \mathbf{h}_i}$.
Since $\frac{\partial \mathbf{c}_t}{\partial \mathbf{h}_i} = \alpha_{t, i} I$, the direct linear term provides constant temporal gradient propagation.
\end{proof}

\section{Complexity Analysis}

\begin{itemize}
    \item \textbf{Time Complexity:} $\mathcal{O}(T_x \cdot (d_a d_s + d_a d_h))$ FLOPs per decoder step.
    \item \textbf{Space Complexity:} $\mathcal{O}(T_x)$ storage for attention weights $\boldsymbol{\alpha}_t$.
\end{itemize}

\section{Exactness and Error Analysis}

The softmax normalization $\sum_i \alpha_i = 1$ ensures that context vector $\mathbf{c}_t$ strictly resides within the convex hull $\text{conv}(\mathbf{h}_1, \dots, \mathbf{h}_{T_x})$, bounding $\|\mathbf{c}_t\| \le \max_i \|\mathbf{h}_i\|$.

\section{Worked Numerical Example}

Let $T_x = 2, d_h = 1, d_s = 1$. Encoder states $\mathbf{h}_1 = [1.0], \mathbf{h}_2 = [3.0]$.
Decoder query state $\mathbf{s} = [2.0]$.
Using Luong dot-product attention ($e_i = \mathbf{s} \cdot \mathbf{h}_i$):
\begin{itemize}
    \item $e_1 = 2.0 \times 1.0 = 2.0$.
    \item $e_2 = 2.0 \times 3.0 = 6.0$.
    \item Normalization: $e_1 - 6.0 = -4.0, e_2 - 6.0 = 0.0$.
    \item $\exp(-4.0) \approx 0.0183, \exp(0.0) = 1.0 \implies \text{Sum} = 1.0183$.
    \item $\alpha_1 = \frac{0.0183}{1.0183} \approx 0.0180, \quad \alpha_2 = \frac{1.0}{1.0183} \approx 0.9820$.
    \item Context vector $c = 0.0180(1.0) + 0.9820(3.0) = 0.0180 + 2.9460 = 2.9640$.
\end{itemize}

\section{Numerical Considerations and Edge Cases}

\begin{itemize}
    \item \textbf{Masking Padding Tokens:} For padded batched sequences, attention energy of pad tokens is set to $-\infty$ ($e_i \gets -10^9$) prior to softmax.
    \item \textbf{Additive vs Multiplicative:} Luong multiplicative attention is faster (GEMM-friendly), while Bahdanau additive attention performs slightly better for high-dimensional dissimilar representations ($d_s \ne d_h$).
\end{itemize}

\section{References}
\begin{enumerate}
    \item Bahdanau, D., Cho, K., & Bengio, Y. (2014). Neural machine translation by jointly learning to align and translate. \emph{ICLR}.
    \item Luong, M.-T., Pham, H., & Manning, C. D. (2015). Effective approaches to attention-based neural machine translation. \emph{EMNLP}, 1412--1421.
\end{enumerate}

\end{document}
